\section{Discussion and Limitations}

\paragraph{What is established.}
Two streams can share the same one-time covariance and the same number of raw
observations yet have different exponents. The decisive object in our model is
the innovation spectrum $\lambda_j/\tau_j$. Uniform persistence produces a
horizontal shift; persistence increasing toward weak modes moves the
model--data frontier. Heavy stationary dwell times add a second, explicitly
quantified phase through residual life.

\paragraph{What is not claimed.}
The theorem does not state that every dependent Gaussian design has exponent
$(b-1)/(a+r)$. The dense AR control shows why: when every sample contains every
mode, temporal correlation need not create an occupancy bottleneck. General
processes may combine innovation coverage, sample-covariance concentration,
and optimization effects. Characterizing that decomposition for arbitrary
noncommuting space--time operators remains open. Representation invariance also
assumes the dense map is known; the ambient minimum-norm result is empirical.
Finally, the real-data proxy estimates $\tau_j$ after PCA and ignores
uncertainty in spectral alignment. A broader benchmark is needed before its
improvement can be treated as universal.

\paragraph{Practical implication.}
A data audit should measure not only variance in direction $j$, but how often
that direction is refreshed. If late spectral directions arrive in long
bursts, more contiguous logging may be less valuable than increasing temporal,
environmental, or population diversity. A global effective sample size can
hide precisely this allocation decision.
